\section{Composition Recipes}
\label{sec:recipes}

A new report starts by copying \texttt{sample/} and editing three blocks. Everything else in
this file can be deleted without breaking the class.

\subsection{Front matter}

The document class line and the branding hooks are the only places where organisation-level
information appears. A released report sets them as follows.

\begin{verbatim}
\documentclass[11pt, a4paper, logo, copyright]{boogonreport}

\renewcommand{\reportorg}{Cooper Intelligence}
\renewcommand{\reportshorttitle}{Cooper Intelligence Technical Report}
\renewcommand{\reportlogo}{../logo/cooper-logo.pdf}

\title{\centering Your Report Title}

\author[*]{Cooper Intelligence\\
  \small under the guidance of Boogon\\
  \small \texttt{official@boogon.com}}

\keywords{three to six terms, semicolon separated}

\begin{abstract}
  One paragraph. It states the question, the method and the result,
  in that order, and stops.
\end{abstract}
\end{verbatim}

\subsection{Sections as separate files}

Long reports keep each section in its own file so that two authors never edit the same file.
The main file then only sequences them, which also makes the order of sections a one-line
decision during review.

\begin{verbatim}
\input{sections/01-purpose}
\input{sections/02-anatomy}
\input{sections/03-typography}
\end{verbatim}

\subsection{The contents page}

Reports of more than a few pages carry a contents page, as the upstream report this class derives
from does. The class provides one command for it: \verb|\reporttoc|, used directly after
\verb|\maketitle|. It starts a new page, sets the contents at 0.9 line spacing so that a long
report still fits them on one page, and starts another new page afterwards.

\begin{verbatim}
\maketitle
\reporttoc
\end{verbatim}

The contents list sections and subsections; deeper levels are excluded by default, which keeps
the page scannable. Page numbers and hyperlinks are generated by \LaTeX{}, so the document is
built twice, as usual, before the numbers settle.

\subsection{Citations and bibliography}

Citations use \texttt{natbib} with author--year output, so that the argument can name a source
in running text. The bibliography style is \texttt{plainnat}; references are kept in
\texttt{refs.bib} next to the main file.

\begin{verbatim}
\citep{lamport1994}    % (Lamport, 1994)
\citet{mittelbach2023} % Mittelbach and Goossens (2023)

\bibliographystyle{plainnat}
\bibliography{refs}
\end{verbatim}

\subsection{Building}

The document is built with \texttt{latexmk}, which runs the required number of \LaTeX{}
passes and invokes \textsc{BibTeX} when the bibliography changes. A clean build from the
repository root is one command.

\begin{verbatim}
latexmk -pdf sample/main.tex
\end{verbatim}

Build artefacts (\texttt{.aux}, \texttt{.log}, \texttt{.pdf}) are ignored by the repository;
only sources are committed. A report is shipped as the pair of PDF and sources, never as a PDF
alone.
